% !TeX root = ../main.tex
\begin{frame}{仿真验证}
  \begin{SlideCanvas}
    \BodyAt{55}{84}{850}{先在单胞和短段圆管中检验变形与承载，再根据主要结果增加血管接触，控制初步研究规模。}
    \PointAt{55}{158}{250}{单胞}
    \PointAt{367}{158}{250}{短段圆管}
    \PointAt{680}{158}{224}{血管接触}
    \AntiCellSketch{113}{220}{.55}
    \MeshTube{381}{225}{190}{79}{BICNavy}
    \foreach \xx in {413,477,541}{
      \draw[-{Latex[length=6pt]},BICDeepRed,line width=1.5pt] (\xx,219)--(\xx,199);
      \draw[-{Latex[length=6pt]},BICDeepRed,line width=1.5pt] (\xx,310)--(\xx,330);}
    \draw[BICNavy,line width=2pt] (690,207)--(890,207) (690,324)--(890,324);
    \draw[BICNavy!40,line width=7pt] (690,196)--(890,196) (690,335)--(890,335);
    \MeshTube{704}{226}{172}{80}{BICAccentBlue}
    \draw[-{Latex[length=8pt]},BICNavy,line width=1.6pt] (292,264)--(341,264);
    \draw[-{Latex[length=8pt]},BICNavy,line width=1.6pt] (604,264)--(653,264);
    \DetailAt{55}{361}{260}{周向拉伸、轴向自由\\
      检查泊松比与曲线误差}
    \DetailAt{367}{361}{260}{径向扩张及规定压缩\\
      比较长度误差、力曲线与局部应变}
    \DetailAt{680}{361}{225}{匹配管腔恢复程度\\
      比较贴壁及壁应变}
    \DetailAt{55}{443}{850}{器件保留轴向自由度；记录端部位移与扭转，并检查网格、加载步长及局部应变。}
    \Sources{拟议验证路径；图形为模型关系示意。超弹性机制验证与永久展开、回弹评价需采用相应本构及加载历史。}
  \end{SlideCanvas}
  \BICConclusionBand{先回答\BICBandEmph{长度与承载能否兼顾}，再评价血管相互作用}
  \note{本页仅展示未来验证设计。第一层沿用课题组超弹性机制模型，第二层保留轴向自由度，仅去除必要刚体运动。长度指标取直径区间内最大绝对相对偏差E_L=max|L(D)/L0-1|；承载拟合误差取同一加载分支上以参考力归一的均方根E_F；另记录局部最大材料应变和端部位移。对目标误差及材料量作参数扫描，形成可行域，而不是只挑选单个成功解。若讨论球囊扩张金属支架永久构型需弹塑性，自膨胀NiTi需相变超弹性及压握释放路径。均匀径向位移用于整体曲线；狗骨效应需允许不同轴向位置独立展开的加载和接触。血管接触在匹配管腔恢复程度后比较间隙、壁应变分布及稳健统计。}
\end{frame}
